Physics-based methods, most notably Density Functional Theory (DFT), provide high accuracy for local bonding environments and band structure prediction. However, the computational scaling of standard Kohn-Sham DFT---typically $\mathcal{O}(N^3)$ with respect to system size---renders large or disordered interface simulations intractable without approximation. Even when feasible, DFT struggles to capture extended features such as long-range strain relaxation, misfit dislocations, and spatially diffuse defect states that dominate realistic interface behaviour \cite{Maji2021-kc, Gaunt1928-kp, Hepplestone2006-lo, Lorentzon2025-si, Low2025-on, Kao2025-ay, Nykiel2023-jf, Gerber2023-yd, Rosen2022-pw, Butler2019-ws, Di_Liberto2022-gz}.

Machine learning potentials (MLPs), such as those based on graph neural networks or equivariant message passing architectures, have emerged as promising tools to extend predictive power to larger systems at near-DFT accuracy \cite{Batzner2022-fd, Lorentzon2025-si, Low2025-on, Batatia2022-xz, Uhrin2025-ba, Owen2024-wp, Rosen2022-pw, Butler2019-ws, Di_Liberto2022-gz}. Nonetheless, they too face structural limitations. MLPs require extensive, representative training datasets, typically derived from DFT calculations, and their performance deteriorates when applied to configurations involving unseen chemistries, surface reconstructions, or high strain regimes \cite{Schutt2019-sz, Leitherer2023-ew, Batzner2022-fd, Lorentzon2025-si, Low2025-on, Owen2024-wp, Rosen2022-pw, Butler2019-ws, Di_Liberto2022-gz}. Furthermore, many ML models lack rigorous uncertainty quantification, making it difficult to detect failure modes in out-of-distribution predictions \cite{Maji2021-kc, Leung2020-tv, Butler2019-ws, Di_Liberto2022-gz}.

Both DFT and ML approaches also fall short in capturing emergent interfacial phenomena such as charge transfer, interface-induced dipoles, or the formation of new interfacial phases. These effects are often strongly dependent on the local atomic registry and chemical heterogeneity, and cannot be inferred from the properties of the isolated constituents alone. Notably, simple models like Anderson\rqss rule have been shown to systematically fail in 2D and 2D|3D heterostructures, necessitating corrections such as $\Delta E_\Gamma$ and $\Delta E_\mathrm{IF}$ that explicitly depend on the atomic-scale interface structure and electrostatics \cite{Maji2021-kc, Low2025-on, Gerber2023-yd, Butler2019-ws, Di_Liberto2022-gz}.

\paragraph{In summary} while ML and physics-based models have complementary strengths, neither is sufficient in isolation for the reliable prediction of interfacial behaviour. Their limitations motivate the development of hybrid workflows, such as embedding physical constraints into ML architectures, or coupling MLP-based screening with selective DFT refinement. These approaches aim to achieve a balance between scalability, accuracy, and generalisability---a balance that remains at the forefront of interface modelling challenges \cite{Maji2021-kc, Leung2020-tv, Batzner2022-fd, Lorentzon2025-si, Low2025-on, Kao2025-ay, Gerber2023-yd, Owen2024-wp, Rosen2022-pw, Butler2019-ws, Di_Liberto2022-gz}.

\subsection{Are there data or generalisation issues?}

The predictive performance of machine learning (ML) models for interfacial systems is fundamentally constrained by data availability, quality, and representational fidelity. While ML models have achieved notable success in bulk materials prediction, their application to interface science introduces distinct challenges; chief among them data scarcity, limited transferability, and inadequate descriptor design.

\subsubsection{Data scarcity and imbalance}

In contrast to bulk crystals, where structured repositories like the Materials Project or AFLOW provide extensive coverage, interface-specific datasets remain sparse \cite{Schutt2019-sz, Butler2019-ws, Di_Liberto2022-gz}. This is due to the combinatorial explosion of viable terminations, slip vectors, and orientations across material pairs, especially in mixed-dimensional systems. Existing databases tend to overrepresent a narrow set of lattice-matched, technologically prominent interfaces, introducing bias and hindering the model's capacity to generalise across broader chemical or structural classes \cite{Maji2021-kc, Low2025-on, Butler2019-ws, Di_Liberto2022-gz}.

Another critical challenge is the lack of standardisation in how interface data is represented. Unlike bulk materials, which can be described relatively easily by a set of structural features such as lattice constants and atomic compositions, interfaces require more complex descriptors, including local bond environments, atomic coordination numbers, and structural motifs that change drastically depending on orientation, termination, and strain. As a result, the design of robust and transferable machine learning models for interfaces is further complicated by the difficulty in developing universally applicable features that capture the intricate nature of interfacial systems \cite{Maji2021-kc, Low2025-on, Butler2019-ws, Di_Liberto2022-gz}. This problem is compounded by the lack of large, consistent datasets that integrate structural, energetic, and electronic data, making it harder for models to learn the subtle interactions that govern interface stability and properties.

Finally, the sheer diversity of possible interface configurations, especially for systems involving two-dimensional (2D) materials, complex oxides, or disordered solids, makes high-throughput DFT calculations or other first-principles methods impractical for exhaustive data collection. To address these issues, there is a pressing need for hybrid data collection strategies that combine high-fidelity, physics-based methods with more scalable, data-driven approaches, such as machine learning potential models that can rapidly generate large datasets of interfacial configurations \cite{Maji2021-kc, Lorentzon2025-si, Low2025-on, Kao2025-ay, Nykiel2023-jf, Batatia2022-xz, Uhrin2025-ba, Gerber2023-yd, Owen2024-wp, Rosen2022-pw, Butler2019-ws, Di_Liberto2022-gz}. Additionally, targeted experimental datasets that validate theoretical predictions are essential for creating reliable and comprehensive interfaces databases.

\subsubsection{Transferability and chemical diversity}

Machine-learned potentials (MLPs) such as MACE \cite{Batatia2022-xz, Uhrin2025-ba, Di_Liberto2022-gz} and ChgNet often fail to extrapolate to configurations that deviate from the training distribution; particularly at strained, chemically heterogeneous, or electronically reconstructed interfaces \cite{Schutt2019-sz, Butler2019-ws, Di_Liberto2022-gz}. The emergence of site-specific hybridisation, localised charge transfer, and dipole formation presents physics not typically encountered in bulk datasets. Without explicit exposure to such environments during training, even state-of-the-art equivariant models exhibit elevated uncertainty and reduced fidelity in these regimes \cite{Khot2025-hh, Maji2021-kc, Leung2020-tv, Fadaly2020-rl, Batzner2022-fd, Low2025-on, Owen2024-wp, Butler2019-ws, Di_Liberto2022-gz}.

\subsubsection{Descriptor and representation limitations}

Generalisation error is further amplified by suboptimal representation of interface-specific physics. Many descriptor schemes, such as SOAP vectors or radial symmetry functions, are optimised for bulk periodicity and fail to capture interfacial asymmetry, vacuum discontinuities, or broken symmetry \cite{Schutt2019-sz, Butler2019-ws, Di_Liberto2022-gz}. Representations that do not enforce scale-invariance or rotational equivariance can also hinder model transferability to supercells or faceted configurations. Physically meaningful descriptors, such as interfacial dipoles, local charge density, or cleavage ene.g. remain underused, though they show promise for enhancing predictive robustness \cite{Maji2021-kc, Low2025-on, Butler2019-ws, Di_Liberto2022-gz}.

\subsubsection{Label quality and training signal noise}

Finally, inconsistencies in training labels present an often-overlooked source of predictive error. Energies obtained from high-throughput DFT pipelines can vary depending on pseudopotential choice, relaxation tolerances, or convergence thresholds. In interfacial contexts, these uncertainties are magnified by metastability and complex relaxation pathways. Some datasets rely on static, unrelaxed geometries or surrogate models, further undermining the reliability of ground truth energetics \cite{Maji2021-kc, Lorentzon2025-si, Low2025-on, Butler2019-ws, Di_Liberto2022-gz}.

\paragraph{In summary} while ML offers substantial speed-ups over DFT, its utility in interface prediction is presently constrained by data bottlenecks, representational mismatch, and transferability failures. Addressing these limitations will require a multi-pronged strategy involving broader and more diverse training datasets, improved descriptors attuned to interfacial physics, and integration of uncertainty quantification into ML workflows. Hybrid schemes, where MLPs perform large-scale screening followed by selective DFT refinement, may offer a viable path forward \cite{Maji2021-kc, Batzner2022-fd, Lorentzon2025-si, Low2025-on, Kao2025-ay, Uhrin2025-ba, Gerber2023-yd, Owen2024-wp, Rosen2022-pw, Butler2019-ws, Di_Liberto2022-gz}.